\documentclass[../main]{subfiles}
\begin{document}

\begin{proposition}
For any \ac{CQ} with nested \texttt{SERVICE} clauses (e.g., \texttt{SERVICE :a \{ ?s ?p ?o . SERVICE :b \{ ?s ?p2 ?o2 \} \}}), provided each intermediate endpoint supports SPARQL~1.1 Federated Query execution, there exists an equivalent \ac{CQ} where all \texttt{SERVICE} clauses are flattened at the \ac{BGP} level (e.g., \texttt{SERVICE :a \{ ?s ?p ?o \} SERVICE :b \{ ?s ?p2 ?o2 \}}).
\end{proposition}

\begin{proof}
We use the FedQPL notation~\cite{Cheng2021}, where $\freq_s^P$ denotes the evaluation of \ac{GP} $P$ at service $s$.
Let $\mathit{tp}_a$ and $\mathit{tp}_b$ be the \ac{BGP} of services $a$ and $b$ respectively.
The nested form corresponds to the following FedQPL expression, where we abuse notation by placing a FedQPL expression in the superscript of $\freq$:
\[
\freq_a^{\fmj(\mathit{tp}_a, \freq_b^{\mathit{tp}_b})}
\]
By assumption, service $a$ supports SPARQL~1.1 Federated Query execution and applies the \ac{GGP} evaluation algorithm of \Cref{List:pseudoCodeFederatedQuery} to its body: it evaluates $\mathit{tp}_a$ locally, obtaining $\freq_a^{\mathit{tp}_a}$, then forwards $\mathit{tp}_b$ to service $b$ as an independent request, obtaining $\freq_b^{\mathit{tp}_b}$.
The two results are joined, giving:
\[
\fmj(\freq_a^{\mathit{tp}_a}, \freq_b^{\mathit{tp}_b})
\]
which is the FedQPL expression for the flattened form.
Since $\fmj$ is commutative and associative and variables have global scope in SPARQL\footnote{\url{https://www.w3.org/TR/2010/WD-sparql11-query-20101014/\#QSynVariables}}, flattening a nested \texttt{SERVICE} clause amounts to reordering operands in the $\fmj$ expression, which does not affect the result.
This reasoning generalizes to any nesting depth.
\end{proof}

\end{document}
